\section{Неки основни примери}
Овде дајемо неке примере из рачунарства који су довољно мали али илустративни да се покакжу неки концепти ТЛА
спецификације.

\subsection{Прекидач за светло}
\begin{primer}[Прекидач за светло]
Прекидач за светло је класичан пример коначног аутомата са два стања: \texttt{off} и \texttt{on}.
Једина акција је притискање прекидача (\texttt{Toggle}) која мења стање из \texttt{off} у \texttt{on} и обрнуто.

\textbf{TLA+ модул (\texttt{LightSwitch.tla}):}
\begin{lstlisting}
---- MODULE LightSwitch ----
VARIABLE light

Init == light = "off"

Toggle ==
    \/ /\ light = "off"
       /\ light' = "on"
    \/ /\ light = "on"
       /\ light' = "off"

Next == Toggle

Spec == Init /\ [][Next]_light

TypeInvariant == light \in {"off", "on"}
====
\end{lstlisting}


\textbf{TLC конфигурациони фајл (\texttt{LightSwitch.cfg}):}
\begin{lstlisting}
INIT Init
NEXT Next
INVARIANT TypeInvariant
\end{lstlisting}

Конфигурациони фајл говори TLC model checkeru:
\begin{itemize}
    \item \texttt{INIT} --- која формула описује почетна стања,
    \item \texttt{NEXT} --- која формула описује прелазе стања,
    \item \texttt{INVARIANT} --- које особине треба да важе у \emph{свим} достижним стањима.
\end{itemize}
Овде први пут видимо нови концепт \textbf{INVARIANT} у конфигурационом фајлу. Он нам говори да неки услов
треба да важи у свим стањима наше спецификације и према томе је инваријантан. 
Особина исправности светло је увек у неком од два дозвољена стања, тј. \texttt{TypeInvariant} дефиниција важи у сваком достижном стању.
    
Такође обратимо пажњу на дефиницију 
\begin{lstlisting}
    TypeInvariant == light \in {"off", "on"}
\end{lstlisting}
овде можемо видети такође везу са математиком. Наиме ово је чиста формула теорије скупова која каже променљива \texttt{light}
припада неком скупу могућих вредности. Могли смо записати $x \in \{p, q\}$, идеја програма би имала исти смисао.
\end{primer}

\begin{primer}[Прекидач за светло --- булова варијанта]
Исти прекидач можемо моделовати и са буловом променљивом уместо стрингова.
Вредност \texttt{TRUE} представља укључено светло, \texttt{FALSE} искључено.
Акција \texttt{Toggle} тада постаје знатно краћа јер нема потребе за две гране.

\textbf{TLA+ модул (\texttt{LightSwitchBool.tla}):}
\begin{lstlisting}
---- MODULE LightSwitchBool ----
VARIABLE light

Init == light = FALSE

Toggle == light' = ~light

Next == Toggle

Spec == Init /\ [][Next]_light

TypeInvariant == light \in BOOLEAN
====
\end{lstlisting}

Уочимо неколико ствари. \texttt{Toggle} је сада само \texttt{light' = \textasciitilde light} негација тренутне вредности.
\texttt{TypeInvariant} користи уграђени скуп \texttt{BOOLEAN} $= \{\texttt{TRUE}, \texttt{FALSE}\}$ уместо ручно набројаног скупа стрингова.
Обе варијанте су семантички еквивалентне; булова је компактнија, а стринг варијанта је читљивија када стања имају смислена имена.
\end{primer}

\subsection{Брава на вратима}
\begin{primer}[Брава на вратима]
Брава на вратима је пример са три стања: \texttt{locked}, \texttt{unlocked} и \texttt{open}.
За разлику од прекидача, овде \emph{нису све акције увек дозвољене} нпр.\ врата се не могу отворити ако су закључана.

Акције су:
\begin{itemize}
    \item \texttt{Unlock} --- откључавање (само из стања \texttt{locked}),
    \item \texttt{Lock} --- закључавање (само из стања \texttt{unlocked}),
    \item \texttt{Open} --- отварање (само из стања \texttt{unlocked}),
    \item \texttt{Close} --- затварање (само из стања \texttt{open}).
\end{itemize}

\textbf{TLA+ модул (\texttt{DoorLock.tla}):}
\begin{lstlisting}
---- MODULE DoorLock ----
VARIABLE door

Init == door = "locked"

Unlock == /\ door = "locked"
          /\ door' = "unlocked"

Lock   == /\ door = "unlocked"
          /\ door' = "locked"

Open   == /\ door = "unlocked"
          /\ door' = "open"

Close  == /\ door = "open"
          /\ door' = "unlocked"

Next == Unlock \/ Lock \/ Open \/ Close

Spec == Init /\ [][Next]_door

TypeInvariant == door \in {"locked", "unlocked", "open"}

CannotOpenWhenLocked == door = "locked" => door # "open"
====
\end{lstlisting}

\textbf{TLC конфигурациони фајл (\texttt{DoorLock.cfg}):}
\begin{lstlisting}
INIT Init
NEXT Next
INVARIANT TypeInvariant
INVARIANT CannotOpenWhenLocked
\end{lstlisting}

Кључна разлика у односу на прекидач је да свака акција садржи \textbf{услов изводљивости} (precondition) у првом конјункту.
На пример, \texttt{Open} може да се изврши само ако је \texttt{door = "unlocked"} ако то не важи, акција је блокирана.
TLC ће аутоматски проверити да у ниједном достижном стању \texttt{CannotOpenWhenLocked} не буде нарушена.
\end{primer}

\subsection{Узајамно искључење}
\begin{primer}[Узајамно искључење]
Узајамно искључење (\textit{mutual exclusion}) је класичан проблем у теорији конкурентних система.
Два процеса се такмиче за приступ критичној секцији, делу кода који смем да се извршава само један процес истовремено.
Сваки процес може бити у једном од три стања: \texttt{idle} (не покушава улазак), \texttt{trying} (чека на улазак) и \texttt{critical} (у критичној секцији).

Овај пример је значајан јер уводи два кључна концепта:
\begin{itemize}
    \item моделовање \textbf{вишеструких процеса} --- стање система је функција која сваком процесу додељује његово тренутно стање,
    \item \textbf{safety особина} --- гарантује да два процеса никад нису истовремено у критичној секцији.
\end{itemize}

\textbf{TLA+ модул (\texttt{MutualExclusion.tla}):}
\begin{lstlisting}
---- MODULE MutualExclusion ----
EXTENDS Naturals

CONSTANT N
ASSUME N \in Nat /\ N > 0

Procs == 1..N

VARIABLE state

States == {"idle", "trying", "critical"}

Init == state = [p \in Procs |-> "idle"]

TryEnter(p) ==
    /\ state[p] = "idle"
    /\ state' = [state EXCEPT ![p] = "trying"]

Enter(p) ==
    /\ state[p] = "trying"
    /\ \A q \in Procs : q # p => state[q] # "critical"
    /\ state' = [state EXCEPT ![p] = "critical"]

Exit(p) ==
    /\ state[p] = "critical"
    /\ state' = [state EXCEPT ![p] = "idle"]

Next == \E p \in Procs : TryEnter(p) \/ Enter(p) \/ Exit(p)

Spec == Init /\ [][Next]_state

TypeInvariant == \A p \in Procs : state[p] \in States

MutualExclusion == \A p, q \in Procs : p # q => ~(state[p] = "critical" /\ state[q] = "critical")
====
\end{lstlisting}

\textbf{TLC конфигурациони фајл (\texttt{MutualExclusion.cfg}):}
\begin{lstlisting}
CONSTANTS
    N = 2

INIT Init
NEXT Next
INVARIANT TypeInvariant
INVARIANT MutualExclusion
\end{lstlisting}

Уочимо неколико важних ствари.
Стање система је \textbf{функција} \texttt{state}, дефинисана над скупом процеса \texttt{Procs = 1..N}.
Израз \texttt{state[p]} даје тренутно стање процеса \texttt{p}, а \texttt{[state EXCEPT ![p] = v]} прави нову функцију идентичну \texttt{state} осим за процес \texttt{p}.

Акција \texttt{Enter(p)} садржи услов \texttt{\textbackslash A q \textbackslash in Procs : q \# p => state[q] \# "critical"} --- процес \texttt{p} може ући у критичну секцију само ако ниједан други процес тренутно није у њој.
Ово је сама суштина узајамног искључења директно изражена у спецификацији.

Особина \texttt{MutualExclusion} формално каже: за свака два различита процеса \texttt{p} и \texttt{q}, не може се десити да су оба у стању \texttt{critical} истовремено.
TLC ће ову особину проверити у свим достижним стањима.
\end{primer}

